\documentclass[../../../hsm_tok_q.tex]{subfiles}

\begin{document}
    \setcounter{figure}{0}
    \textsf{図\ref{fig/hsm_tok_2014_2nd_05_01_q}}に示した立体A--BCDEは，%
    底面BCDEが1辺の長さ8\,$\mathrm{cm}$の正方形で，%
    AB＝AC＝AD＝AE＝8\,$\mathrm{cm}$の正四角すいである．

    点Pは，頂点Bを出発し，辺BA，辺AD上を毎秒1\,$\mathrm{cm}$の速さで動き，%
    16秒後に頂点Dに到着する．

    次の各問に答えよ．
    %
    \vspace{.5\baselineskip}
    {\par \centering
        \setlength{\abovecaptionskip}{3pt}
        \includegraphics%
            {\subfix{fig_hsm_tok_2014_2nd_05/fig_hsm_tok_2014_2nd_05_01_q.pdf}}
            \captionof{figure}{} \label{fig/hsm_tok_2014_2nd_05_01_q}
    \par}
    \vspace{.4\baselineskip}
    %
    \begin{enumerate}
        \item 次の\:\myboxa{　　　}\:の中の「\textsf{か}」「\textsf{き}」に当てはまる
            数字をそれぞれ答えよ．

            頂点Dと点Pを結んだ場合を考える．

            点Pが頂点Bを出発してから4秒後のとき，線分DPの長さは，%
            \:\myboxa{　\textsf{か}　}\:$\sqrt{\,\text{\myboxa{　\textsf{き}　}}\,}$\:%
            cmである．
        %
        \item 次の\:\myboxa{　　　}\:の中の「\textsf{く}」「\textsf{け}」「\textsf{こ}」%
            に当てはまる数字をそれぞれ答えよ．

            \textsf{図\ref{fig/hsm_tok_2014_2nd_05_02_q}}は，%
            \textsf{図\ref{fig/hsm_tok_2014_2nd_05_01_q}}において，%
            点Pが頂点Bを出発してから13秒後のとき，点Pから辺DEに垂線を引き，%
            辺DEとの交点をQとし，頂点Bと点P，頂点Bと点Q，頂点Cと点P，頂点Cと点Qを%
            それぞれ結んだ場合を表している．

            立体P--BCQの体積は，
            \:\myboxa{\hspace{0.45\zw}\textsf{く\hspace{0.1\zw}け}\hspace{0.45\zw}}%
            $\sqrt{\,\text{\myboxa{　\textsf{こ}　}}\,}$\:%
            $\mathrm{cm}^3$である．
            \\
            \begin{minipage}{\linewidth}
                \vspace{.5\baselineskip}
                {\par \centering
                    \setlength{\abovecaptionskip}{3pt}
                    \includegraphics%
                        {\subfix{fig_hsm_tok_2014_2nd_05/fig_hsm_tok_2014_2nd_05_02_q.pdf}}
                        \captionof{figure}{} \label{fig/hsm_tok_2014_2nd_05_02_q}
                \par}
            \end{minipage}
    \end{enumerate}
\end{document}


